\id{ҒТАМР 65.63.33}{}

\begin{header}
\swa{С.А.~Карденов, А.К.~Игенбаев, Ш.Б.~Байтукенова, С.Б.~Байтукенова, Ж.С.~Ажгереева}{ГУМИНДІ ЗАТТАРДЫҢ СҮТҚЫШҚЫЛДЫ ӨНІМДЕРДІҢ САПАЛЫҚ КӨРСЕТКІШТЕРІНЕ ЫҚПАЛЫ}

\tsp{1}С.А.~Карденов\envelope,
\tsp{1}А.К.~Игенбаев,
\tsp{1}Ш.Б.~Байтукенова,
\tsp{2}С.Б.~Байтукенова,
\tsp{1}Ж.С.~Ажгереева
\end{header}

\begin{affil}
\tsp{1}«С.Сейфуллин атындағы Қазақ агротехникалық зерттеу университеті» КеАҚ, Астана, Қазақстан,

\tsp{2}«Қ.Құлажанов атындағы Қазақ технология және бизнес университеті» АҚ, Астана, Қазақстан

\corrauthor{Корреспондент-автор: Askerbekovsk@mail.ru}
\end{affil}

Бұл мақалада қазіргі таңдағы өзекті мәселелердің біріне айналған сүт
және сүтқышқылды өнімдер жайында болмақ. Сондай-ақ мақалада отандық және
шетелдік ғалымдардың ғылыми жұмыстары мен зертеулері қарастырылып, соңғы
жылдардағы жүргізілген зерттеу жұмыстарына байланысты жаңа технологиясы
жетілдірілді.

Сүтқышқылды өнімдердің технологиясында гуминді заттарды қолдану
мүмкіндіктері зерттеліп, олардың осы өнімдерге әсер ету ерекшеліктері
көрсетілді. Гуминдік заттар топырақ пен судағы ең көп таралған
органикалық қосылыстардың бірі болып табылады. Олар сүтқышқылды
өнімдердің биологиялық құндылықтарына оң әсер ететін ерекше қасиеттерге
ие. Гуминдік заттар биологиялық белсенді қосылыстар ретінде жаңа
биологиялық белсенді заттардың көзі болып табылады. Калий гуматымен
байытылған сүтқышқылды өнім айран өндіру технологиясын ойдағыдай жүзеге
асыруға болатынын айта аламыз, осыған орай айранды гуматпен өндіру
технологиясын жетілдірдік. Тәжірибелік зерттеулер нәтижесінде майлылығы
2,5\% 1 л айранға 0,5\% калий гуматын қосу негізделді.

Калий гуматымен байытылған айранның бірегей химиялық құрамы оның
коммерциялық әлеуетін арттырады. Айрандағы сүтқышқылды бактериялар мен
ұйытқылар оны пробиотиктерге жа\-қындатады және сақтау кезінде
органолептикалық қасиеттерін жақсартады. Бұл өнім әсіресе дәрумендер мен
микроэлементтердің жетіспеушілігінен зардап шегетін адамдар үшін өте
қажет. Айран ішек микрофлорасына оң әсер етіп, дұрыс тамақтануға ықпал
етеді.

Зерттеу нәтижелері бойынша гумин қышқылдарымен байытылған айран алудың
технологиясы мен рецептурасы жасалды. Бұл технология өндірісте қолдануға
ыңғайлы. Калий гуматының ең тиімді концентрациясы 0,5\% құрайды.
Тәжірибе нәтижесінде гумин қосылған айранды тұтынуда теріс әсерлер мен
асқынулар анықталмады және гуминдік заттармен байытылған айран өнімінің
химиялық құрамын, энергетикалық құндылығын, микробиологиялық және рН
көрсеткіштерін анықтадық.

{\bfseries Түйін сөздер:} гуминдік заттар, калий гуматы, сүтқышқылды
өнімдер, айран, сүтқышқылды бактериялар, "Қазгумат" препараты,
минералды-органикалық тыңайтқыш, дәрумендер.

\begin{header}
ВОЗДЕЙСТВИЕ ГУМИНОВЫХ ВЕЩЕСТВ НА КАЧЕСТВО КИСЛОМОЛОЧНЫХ ПРОДУКТОВ

\tsp{1}С.А.~Карденов\envelope,
\tsp{1}А.К.~Игенбаев,
\tsp{1}Ш.Б.~Байтукенова,
\tsp{2}С.Б.~Байтукенова,
\tsp{1}Ж.С.~Ажгереева
\end{header}

\begin{affil}
\tsp{1}НАО «Казахский агротехнический исследовательский университет им. С.Сейфуллина», Астана, Казахстан,

\tsp{2}АО «Казахский университет технологии и бизнеса им. К. Кулажанова», Астана, Казахстан,

e-mail: Askerbekovsk@mail.ru
\end{affil}

В данной статье рассматриваются молоко и молочные продукты, ставшие
одной из самых актуальных тем современности. Также дан обзор научных
трудов и исследований отечественных и зарубежных ученых, а также новые
технологии, разработанные и усовершенствованные в последние годы
благодаря исследовательской работе.

В статье изучены возможности использования гуминовых веществ в
технологии кисломолочных продуктов и показаны особенности их воздействия
на эти продукты. Гуминовые вещества являются одними из наиболее
распространенных органических соединений в почве и воде. Они обладают
уникальными свойствами, которые положительно влияют на биологическую
ценность кисломолочных продуктов. Гуминовые вещества являются источником
новых биологически активных соединений. Можно сказать, что технология
производства кефира - молочнокислого продукта, обогащенного гуматом
калия, может быть успешно реализована, в связи с этим нами
усовершенствована технология производства кефира с гуматом калия.
Экспериментальными исследованиями обоснована целесообразность добавления
0,5\% гумата калия на 1 л кефира жирностью 2,5\%.

Уникальный химический состав кефира, обогащенного гуматом калия,
повышает его коммерческий потенциал. Молочнокислые бактерии и дрожжи в
кефире приближают его к пробиотикам и улучшают его органолептические
свойства при хранении. Этот продукт особенно необходим людям, страдающим
от недостатка витаминов и микроэлементов. Кефир положительно влияет на
микрофлору кишечника и способствует правильному питанию.

На основании результатов исследований разработаны технология и рецептура
производства кефира, обогащенного гуминовыми кислотами. Эту технологию
удобно использовать в производстве. Наиболее эффективная концентрация
гумата калия -- 0,5\%. В результате эксперимента не выявлено
отрицательных эффектов и осложнений от употребления кефира с гуминовыми
веществами, определены химический состав, энергетическая ценность,
микробиологические и рН показатели кефирного продукта, обогащенного
гуминовыми веществами.

{\bfseries Ключевые слова:} гуминовые вещества, гумат калия, кисломолочные
продукты, кефир, кисломолочные бактерии, препарат «Казгумат»,
минерально-органическое удобрение, витамины.

\begin{header}
EFFECT OF HUMIC SUBSTANCES ON THE QUALITY OF FERMENTED MILK PRODUCTS

\tsp{1}S.A.~Kardenov\envelope,
\tsp{1}A.K.~Igenbayev,
\tsp{1}Sh.B.~Baitukenova,
\tsp{2}S.B.~Baitukenova,
\tsp{1}Zh.S.~Azhgereeva
\end{header}

\begin{affil}
\tsp{1}NAO «S. Seifullin Kazakh Agrotechnical Research University», Astana, Kazakhstan,

\tsp{2}JSC «K. Kulazhanov Kazakh University of Technology and Business» Astana, Kazakhstan,

e-mail: Askerbekovsk@mail.ru
\end{affil}

This article discusses milk and dairy products, which have become one of
the most pressing issues of our time. It also provides an overview of
scientific works and research by domestic and foreign scientists, as
well as new technologies developed and improved in recent years thanks
to research work.

The article examines the possibilities of using humic substances in the
technology of fermented milk products and shows the features of their
impact on these products. Humic substances are among the most common
organic compounds in soil and water. They have unique properties that
have a positive effect on the biological value of fermented milk
products. Humic substances are a source of new biologically active
compounds. It can be said that the technology of producing kefir - a
fermented milk product enriched with potassium humate, can be
successfully implemented, and in this regard, we have improved the
technology of producing kefir with potassium humate. Experimental
studies substantiate the feasibility of adding 0.5\% potassium humate
per 1 liter of kefir with a fat content of 2.5\%.

The unique chemical composition of kefir enriched with potassium humate
increases its commercial potential. Lactic acid bacteria and yeast in
kefir bring it closer to probiotics and improve its organoleptic
properties during storage. This product is especially necessary for
people suffering from a lack of vitamins and microelements. Yogurt has a
positive effect on the intestinal microflora and promotes proper
nutrition.

Based on the research results, a technology and recipe for the
production of kefir enriched with humic acids have been developed. This
technology is convenient to use in production, saving time and money.
The most effective concentration of potassium humate is 0.5\%. As a
result of the experiment, no negative effects and complications from the
use of kefir with humic substances were revealed, the chemical
compo\-sition, energy value, microbiological and pH indicators of the
kefir product enriched with humic substances were determined.

{\bfseries Keywords:} Humic substances, potassium humate, lactic acid
products, kefir, lactic acid bacteria, the drug "Kazgumat",
mineral-organic fertilizer, vitamins.

\begin{multicols}{2}
{\bfseries Кіріспе.} Қазіргі тамақ өнеркәсібі негізгі азықтық
қажеттіліктерді қанағаттандырумен қатар адам денсаулығына оң әсерін
тигізетін, биологиялық белесді компоненттермен байытылған функционалды
өнімдерге деген сұраныстың өсуімен сипатталады. Бұл контексте гумин
заттары (ГЗ) - өсімдік және жануар қалдықтарының ыдырауы нәтижесінде
түзілетін табиғи жоғары молекулалық полимерлер - ерекше қызығушылық
тудырады {[}1{]}. Биологиялық белесінің кең спектріне (қабынуға қарсы,
антиоксиданттық, сорбциялық) ие болуына байланысты, ГЗ азық өнімдерін
байыту үшін перспективалы ингредиенттер болып табылады. Алайда оларды
сүт өнеркәсібінде, әсіресе, қышқыл сүт өнімдерінде қолдану бірқатар
технологиялық және органолептикалық қиындықтармен байланысты. Осылайша,
бұл зерттеудің мақсаты гумин заттарының қышқыл сүт өнімдерінің сапасына
қатысты физика-химиялық, микробиологиялық және органолептикалық
көрсеткіштеріне әсерін кешенді талдау болып табылады. Зерттеу гипотезасы
ГЗ-тың оңтайландырылған дозалары ферментация процесіне және өнімнің
сезгіш сипаттамаларына жағымсыз әсерін тигізбей, өнімнің функционалдық
қасиеттерін жақсарта алады деген болжамға негізделген.

Әлеуметтік -- экологиялық проблемалар. Жер шарындағы экологиялық
жағдайдың нашарлауы, халықтың дәрумендер мен микроэлементтердің
жетіспеушілігінен зардап шегуі, қала ауасының ластануы және патогенді
микроорганизмдердің таралуы өзекті мәселелерге айналды. Гуминды заттар
табиғи органикалық заттардың және қоршаған ортада кездесетін ең көп
таралған органикалық қосылыстардың өкілі ретінде қоршаған ортаны тазарту
және қалпына келтіру үшін қолданылады {[}2{]}. Дәрумендер - адамдардың
денсаулығын сақтауда шешуші рөл атқарады. Олардың тотығу, қабыну,
иммундық, генетикалық және эпигенетикалық реттеудегі рөлі оларды
жасушалық гомеостазды сақтауды үшін мағызды етеді {[}3{]}.

Гумин заттары - топырақта, шөгінділерде және су айдындарында кең
таралған, топырақ органикалық заттарының негізгі құрамдас бөлігін 60-70
құрайтын табиғи макромолекулалық органикалық қосылыс {[}4{]}.

Гумин қышқылдары -- топырақ қасиеттерің, өсімдіктердің өсуін және
агрономиялық көрсеткіштерді жаөсартуда маңызды рөл атқаратын органикалық
молекулалар. Гуминды қышқылдар көздеріне көмір, қоңыр көмір, топырақ
және органикалық заттар жатады. Соңғы жылдары ауыл шаруашылығы
өндірісінің тұрақтылығын қамтамасыз ету үшін дақылдарды өндіруде гумин
қышқылына негізделген өнімдер қолданылуда {[}5{]}.

Дәл осы функционалдық топтар олардың негізгі қасиеттерін анықтайды:

1. Антиоксиданттық қасиеттер - металл иондарын хелаттау және еркін
радикалдарды бейтараптау қабілеті гумин заттарын табиғи тиімді
антиоксидантқа айналдырады, бұл өнімдердің сақталу мерзімін ұзартуға
мүмкіндік береді.

2. Пребиотикалық әлеует - гумин заттарының кейбір фракциялары, атап
айтқанда төмен молекулалық фракциялардағы гумин қышқылдары, ішек
микрофлорасының пайдалы өкілдері (бифидобактериялар және
лактобациллалар сияқты) үшін селективті субстрат ретінде қызмет ете
алады.

3. Антимикробтік қасиеттер - гумин заттары бірқатар патогенді және шартты
түрде патогенді микроорганизмдерге (мысалы, Escherichia coli,
Staphylococcus aureus) қарсы тежеуші әсер көрсетеді, бірақ сүтқышқылды
ашытқы саңырауқұлақтарының белесін бәсеңдетпейді.

4. Сорбциялық қабілет - гумин заттары улы заттарды, ауыр металдарды және
аллергендерді байланыстырып, шығаруға қабілетті, бұл өнімнің
детоксикациялық қасиеттерін күшейтеді.

Гуминдік заттар топырақ пен судағы ең көп таралған органикалық
қосылыстардың бірі болып табылады. Олар сүтқышқылды өнімдердің
биологиялық құндылықтарына оң әсер ететін ерекше қасиеттерге ие {[}6{]}.

Сүтқышқылды өнімдер елдің мәдени мұрасында терең тамыр жайған бұл сүт
өнімдері бай тағамдық құндылығы үшін ұзақ уақыт бойы бағаланып келеді.
Дәмінен басқа, бұл сүтқышқылды өнімдер пробиотикалық микроорганиздердің
көзі ретінде айтарлықтай әлеуетке ие, бұл тұтынушылардың денсаулығына
пайдалы болуы мүмкін {[}7{]}.

Гуминдік заттар биологиялық белсенді қосылыстар ретінде жаңа биологиялық
белсенді заттардың көзі болып табылады. Мысалы, медициналық лигнин
(полифепан) клиникалық тәжірибеде кеңінен қолданылады. Гуматтар -
өсімдік талшықтарының ыдырауы кезінде топырақта түзілетін табиғи
элементтер, натрий немесе калий тұздары түрінде болады {[}1{]}.

Гумин қышқылдары негізінде жасалған отандық препараттардың
биорегуляциялық қасиеттері микробқа қарсы белсенділігін, биологиялық
белсенділігін және зиянсыздығын көрсетті. Бұл гумин қышқылдарының
экологиялық қауіпсіз отандық препараттарын қолданудың болашағы зор
екенін көрсетеді.

Зерттеудің өзектілігі - гуминдік заттардың сүтқышқылды өнімдердің
биологиялық құндылығын зерттеу және жетілдіру. Жаңалығы -- гуминдік
заттардың сүтқышқылды өнімдердің технологияларында қолдану мүмкіндігі,
әсер ету ерекшеліктері және гуминдық заттармен байытылған дайын өнімнің
функционалдық-технологиялық және құрылымдық-механикалық қасиеттерінің
өзгеру сипаты зерттеледі.

Ғылыми зерттеудің мақсаты - Қазақстанда өндірілетін көмірді, әсіресе
Сарыадыр кен орнының көмірін тотықтыру арқылы «Көмір химиясы және
технологиясы институты» ЖШС қазақстандық ғалымдары алған «Қазгумат»
препаратынан калий гуматы негізінде айран алудың рецептурасы мен
технологиясын жасау.

Осы мақсатқа жету үшін келесі міндеттер қойылды:

- гуминдік заттарды алу тәсілдерін таңдау және негіздеу;

- гуминдік заттардың сүтқышқыл өнімі сынамасындағы өзгерістерін және
сапасына әсерін зерттеу;

- гуминдік заттарды ғылыми негіздеу және физика-химиялық қасиеттерін
зерттеу нәтижесінде айран өндіру технологиясында тиімді қолдану арқылы
функционалды бағыттағы өнім алу;

- айран өнімінің сапалық көрсеткіштерін және гуминдік заттардың әсерін
зерттеу;

- гуминдік заттармен байытылған айран өнімінің химиялық құрамын,
энергетикалық құндылығын, микробиологиялық және рН көрсеткіштерін
анықтау.

{\bfseries Материалдар мен әдістер.} Зерттеу объектісі ретінде «Қазгумат»
препараты негізіндегі калий гуматымен (0,1\%; 0,5\% және 1\%
концентрацияда) байытылған айран таңдалды. Бақылау үлгісі ретінде
қоспасыз дәстүрлі әдіспен дайындалған айран қолданылды.

Негізгі зерттеу материалдары: майлылығы 2,5\% болған жаңа сауылған сиыр
сүті, сапасы және қауіпсіздігі қазақстандық стандарттарға сәйкес келеді.

Дәстүрлі айран алу үшін қолданылатын бактериялық ашытқы~Kefir 30 ұйытқы
сынамы (әдістемеде көрсетілгендей 10 мл).

Байытқыш~«Қазгумат» препараты негізіндегі калий гуматы. Препарат «Көмір
химиясы және технологиясы институты» ЖШС ғалымдарының Сарыадыр кен
орнының В2 маркалы қоңыр көмірін тотықтыру арқылы алынған. Бұл әдіс
ультрадыбыстық өңдеуді және атмосфералық оттегіні қолдануды қамтиды,
нәтижесінде суда еритін металлорганикалық хелаттық кешендері бар калий
гуматы түзіледі.

Айран өндірісінің технологиясы - калий гуматымен байытылған айран
С.Сейфуллин атындағы Қазақ агротехникалық зерттеу университетінің
тәжірибелік-өндірістік цехында резервуарлық әдіспен дайындалды.
Технологиялық процесс 1-сұлбаға сәйкес жүргізілді.

Зерттеу жұмысы кешенді түрде жүргізіліп, өнімнің сапасы мен
қауіпсіздігінің келесі көрсеткіштері бағаланды: Қышқылдық титрлеу
әдісімен (°Т) анықталды. рН көрсеткіші рН-метр көмегімен өлшенді.
Майлылық гравиметриялық әдіспен (Гербер әдісі) анықталды. Ақуыз құрамы
жоғары өнімді сұйық хроматография әдісімен анықталды. Макро- және
микроэлементтер құрамы атом-абсорбциялық спектрометрия және
масс-спектрометрия әдістері қолданылып, кальций, темір және басқа
элементтердің мөлшері анықталды.

Калий гуматының сипаттамасы: «Қазгумат» препаратының элементтік құрамы
(С, Н, О, N, S), функционалдық топтары (карбоксил, фенолды) және
органолептикалық көрсеткіштері (түс, рН) сәйкес мемлекеттік стандарттар
бойынша зерттелді.

Микробиологиялық талдаулар жүргізу барысында микроорганизмдердің жалпы
саны стандартты агарлық орталардағы есу әдісімен анықталды.

Органолептикалық бағалау параметрлері арқылы анықталды: түс (ақтан сүт
түстес аққа дейін), консистенция (біртектілік, тұтқырлық), иіс (таза,
сүтқышқылды, бөтен иіссіз) және дәм (таза, ашықшыл, сәл ашқылтым).

Әрбір сынама (0,1\%, 0,5\%, 1\% гумат қосылған және бақылау үлгісі) 5
балдық жүйе бойынша бағаланды.

Алынған барлық нәтижелер статистикалық өңдеуден өткізілді. Әрбір
тәжірибе 3-5 рет қайталанды. Алынған нәтижелерді өңдеу үшін вариациялық
статистика әдістері (орташа мән, стандартты ауытқу) қолданылды.
Нәтижелердің статистикалық маңыздылығы Стьюдент критериі бойынша
(p<0,05 деңгейінде) бағаланды.

{\bfseries Нәтижелер және талқылау.} Ферментация үдерісі мен
микробиологиялық көрсеткіштерге әсері. Зерттеулер гумин заттарының (ГЗ)
сүтқышқылды микроорганизмдеріне әсерінің екіжақты екенін көрсетеді.
Төмен мөлшерде (0,001-0,01\%) ГЗ пробиотикалық сипаттамалары
(Lactobacillus acidophilus, Bifidobacterium spp.) өсуінің ынталандырғышы
ретінде әрекет ете алады. Бұл олардың жасуша мембраналарының қоректік
заттарға өткізгіштігін жақсарту және ортаның рН-ын тұрақтандыра отырып,
буферлік функцияны атқару қабілетімен байланысты. Алайда, мөлшері
жоғарылаған кезде (\textgreater0,05\%) қышқыл түзілу жылдамдығының
баяулауы байқалуы мүмкін, бұл бактериялар метаболизмі үшін қажет кальций
иондарын байланыстырумен және гумин заттарының күшті антимикробтік
әсерімен түсіндіріледі {[}8{]}.

Сондықтан, негізгі параметр - гумин заттарының ашытқы үдерісін
тежеместен, пребиотикалық әсері қамтамасыз ететін мөлшерін таңдау болып
табылады.

Физика-химиялық және құрылымды-механикалық қасиеттерге әсері. Гумин
заттары реологиялық сипаттамалар мен ұйытқың құрылымына айтарлықтай әсер
етеді.

Гумин заттарының оңтайлы мөлшерінде олар рН-дың төмендеу динамикасын
бұзбайды, бұл стандартты ашыту уақытын қамтамасыз етеді.

Гумин заттарын 0,005-0,02\% мөлшерінде енгізу айранның тұтқырлығын
арттыруға және сұйықтық бөлінуінің (синерезис) азаюына ықпал етеді. Бұл
гумин заттары сүт ақуыздарымен (казеин) өзара әрекеттесуінің нәтижесінде
ылғалды жақсы ұстайтын берік ақуыздық-гумин матрицасының қалыптасуымен
түсіндіріледі {[}9{]}.

Гумин заттарының антиоксиданттық қасиеттері майлардың тотығуына жол
бермейді, бұл майлылығы жоғары өнімдер (қаймақ, ірімшік) үшін өте
маңызды болып табылады және олардың органолептикалық тұрақтылығын
ұзартады.

Зерттеу материалы ретінде елімізде өндірілетін калий гуматы қолданылды.
Ол «Көмір химиясы және технологиясы институты» ЖШС қазақстандық
ғалымдарының Сарыадыр кен орнының көмірін тотықтыру арқылы алынған
«Қазгумат» препаратынан өндірілді. Содан кейін С.Сейфуллин атындағы
Қазақ агротехникалық зерттеу университетінің сүт өнімдерін қайта өңдеуге
арналған тәжірибелік-өндірістік цехында резервуарлық әдіспен айран
дайындалды. Жаңа сүтқышқылды өнімді алу технологиясы 0,1\%, 0,5\% және
1\% калий гуматының концентрацияларымен байытылды, қышқылды өнім алу
технологиясы 1-ші сұлба көрсетілген.
\end{multicols}

\fig[0.45\textwidth]{p/image51}[Сұлба 1 - Калий гуматымен байытылған айранның технологиялық сұлбасы]

\begin{multicols}{2}
«Қазгумат» - қоңыр көмірді тотықтыру арқылы алынатын қара қоңыр зат -
калий гуматы негізіндегі тұрмыстық минералды-органикалық тыңайтқыш.
Тыңайтқышты алу технологиясын «Көмір химиясы және технологиясы
институты» жауапкершілігі шектеулі серіктестігі мен өндіруші
«Казтехноуголь» ЖШҚ ғылыми-өндірістік бірлестігінің авторлар тобы
әзірледі. Технология ультрадыбыстық және атмосфералық оттегінің тотыққан
B2 маркалы қоңыр көмірге біріктірілген әсеріне негізделген, бұл ұсақ
мицеллалардың пайда болуына әкеледі. Осылайша алынған калий гуматы
құрамында C--O, C--OH және C--NH бос функционалды топтары бар суда
еритін металлорганикалық хелат кешендері бар {[}10{]}. Талдау
нәтижелерінің негізінде қолда бар қоректік заттар 1-кестеде
көрсетілгендей келесі формаларда:
\end{multicols}

\tcap{1-кесте. «Қазгумат» құрамындағы қоректік заттар}
\begin{longtblr}[
  label = none,
  entry = none,
]{
  cells = {c},
  cells = {font = \small},
  row{1} = {font=\bfseries\small},
  hlines,
  vlines,
}
№ & Қоректік заттар & Құрамы, \% \\
1 & Гуминдік заттар & 53-59      \\
2 & Калий (К)       & 12.53      \\
3 & Азот (N)        & 1.1        \\
4 & Оттегі (O)      & 16.57      \\
5 & Кремний (Sі)    & 16.7       
\end{longtblr}

\begin{multicols}{2}
2-ші кестеде «Қазгумат» калий гуматының химиялық құрамы, 3-ші кестеде
«Қазгумат» калий гуматының сипаттамасы, 4-ші кестеде органолептикалық
және 5-ші кестеде рентгенологиялық көрсеткіштері, ал 6-шы кестеде калий
гуматының химиялық құрамы көрсетілген.
\end{multicols}

\tcap{2-кесте. Көміртегінің (С), сутегінің (Н), азоттың (N), күкірттің
(S) және оттегінің (O) оташа мөлшері}
\begin{longtblr}[
  label = none,
  entry = none,
]{
  colspec = {Q[123]Q[69]Q[69]Q[62]Q[69]Q[79]Q[63]Q[63]Q[62]Q[63]Q[204]},
  cells = {c},
  cells = {font = \small},
  row{1} = {font=\bfseries\small},
  cell{1}{1} = {r=2}{},
  cell{1}{2} = {c=5}{},
  cell{1}{7} = {c=4}{},
  cell{3}{1} = {r=3}{},
  cell{3}{11} = {r=2}{},
  cell{4}{2} = {c=5}{},
  cell{4}{7} = {c=4}{},
  hlines,
  vlines,
}
Атауы        & Құрамы, \%    &        &       &        &         & Атомдық қатынас  &       &       &       & Химиялық формуласы \\
             & С             & Н      & N     & S      & О       & С/Н              & С/О   & С/N   & С/S   &                    \\
Калий гуматы & 22.205        & 1.9065 & 3.495 & 0.6775 & 20.4435 & 0.979            & 0.447 & 7.411 & 87.49 & С95Н95О65N15 S1    \\
             & Атомдық үлесі &        &       &        &         & Пайыздық қатынас &       &       &       &                    \\
             & 1.85          & 1.89   & 0.25  & 0.02   & 1.28    & С/Н              & С/О   & С/N   & С/S   &                    
\end{longtblr}

\tcap{3-кесте. «Қазгумат» калий гуматының сипаттамасы}
\begin{longtblr}[
  label = none,
  entry = none,
]{
  width = \linewidth,
  colspec = {Q[71]Q[58]Q[50]Q[63]Q[46]Q[46]Q[42]Q[42]Q[42]Q[42]Q[42]Q[46]Q[42]Q[42]Q[140]Q[98]},
  cells = {c},
  cells = {font = \footnotesize},
  row{1} = {font = \bfseries\footnotesize},
  cell{1}{1} = {r=2}{},
  cell{1}{2} = {c=3}{0.1\linewidth},
  cell{1}{5} = {c=10}{0.5\linewidth},
  cell{1}{15} = {c=2}{0.15\linewidth},
  hlines,
  vlines,
}
Үлгі         & Техникалық сипаттама, \% &     &       & Құрғақ затқа шаққандағы элемент мөлшері, мас. \% &       &      &      &      &      &      &       &      &      & Фуникционалдық топтардың сипаттамасы, мг-экв/г &       \\
             & Wrt                      & Ar  & Vd    & C                                                                     & O     & Na   & Mg   & Al   & Si   & S    & K     & Fe   & F    & Карбоксил                                      & Фенол \\
Калий гуматы & 44.8                     & 5.5 & 42.69 & 35.95                                                                 & 30.16 & 0.44 & 0.25 & 8.57 & 1.12 & 0.14 & 21.05 & 1.47 & 0.84 & 0.5                                            & 0.7   
\end{longtblr}

\tcap{4-кесте. Калий гуматының органолептикалық көрсеткіштері}
\begin{longtblr}[
  label = none,
  entry = none,
]{
  width = \linewidth,
  colspec = {Q[20]Q[260]Q[256]Q[175]Q[204]},
  cells = {c},
  cells = {font = \small},
  cell{1}{2} = {font=\bfseries\small},
  cell{1}{3} = {font=\bfseries\small},
  cell{1}{4} = {font=\bfseries\small},
  cell{1}{5} = {font=\bfseries\small},
  hlines,
  vlines,
}
№ & Көрсеткіш атауы                    & Норма НБ                              & Нақты көрсеткіштер       & Сынау әдістері бойынша НБ \\
1 & Сыртқы түрі, түсі                  & Қара қоңырдан қараға дейінгі сұйықтық & Қою қоңыр түсті сұйықтық & МЕМСТ 9097-82             \\
2 & Сутекті көрсеткіш, рН              & 11.6                                  & 10-11.6                  & МЕМСТ 27979-88            \\
3 & Гумин қышқылдарының массалық үлесі & 54.86                                 & 56                       & МЕМСТ 9517-94             
\end{longtblr}

\begin{multicols}{2}
3 кестеде келтірілген органолептикалық сипаттамаларға әсері көрсетілген.
Бұл ең сыни аспект болып табылады. Гумин заттары тіпті аз мөлшерде
енгізу өнімге тұтынушы үшін әдеттен тыс болатын ашық-қоңыр немесе крем
түс береді. Сонымен қатар, жеңіл «топырақты» немесе «дегть тәрізді» дәмі
пайда болады. Бұл әсерлерді азайту үшін төмендегі қағидаттар
қолданылады:

- түсі мен дәмі азырақ байқалатын тазартылған (ультрасүзгіленген) ГЗ
фракцияларын пайдалану;

- ГЗ-ті бөтен дәмдерді бүркеп тастайтын хош иістендіргіштермен (ваниль,
какао, жидектер) біріктіру;

- ГЗ-ті микрокапсулаланған түрде енгізу, бұл олардың өнімнің дәм-хош иіс
матрицасымен тікелей байланысынан оқшаулайды {[}11{]}.
\end{multicols}

\tcap{5-кесте. Калий гуматының рентгенологиялық көрсеткіші}
\begin{longtblr}[
  label = none,
  entry = none,
]{
  width = \linewidth,
  colspec = {Q[587]Q[85]Q[40]Q[182]},
  cells = {font = \small},
  hlines,
  vlines,
}
Табиғи радионуклидтердің меншікті тиімді белсенділігі, БК/кг & Көп емес & 12.0 & МЕМСТ 301008-94 
\end{longtblr}

\tcap{6-кесте. Калий гуматының химиялық құрамы}
\begin{longtblr}[
  label = none,
  entry = none,
]{
  cells = {c},
  cells = {font = \small},
  row{1} = {font=\bfseries\small},
  row{2} = {font=\bfseries\small},
  cell{1}{1} = {r=2}{},
  cell{1}{2} = {r=2}{},
  hlines,
  vlines,
}
№  & {Анықталатын көрсеткіштер,\\мкг/кг 1х10\textsuperscript{-6}} & Тұтынушы үлгісі \\
   &                                                                                         & Калий гуматы    \\
1  & Алюминий (Al)                                                                           & 1497            \\
2  & Ванадий (V)                                                                             & 4.91            \\
3  & Гадолиний (Gd)                                                                          & 1.80            \\
4  & Диспрозий (Dy)                                                                          & 2.04            \\
5  & Темір (Fe)                                                                              & 4849            \\
6  & Иттербий (Yb)                                                                           & 0.900           \\
7  & Иттрий (Y)                                                                              & 5.33            \\
8  & Калий (K)                                                                               & 28.1            \\
9  & Калций (Ca)                                                                             & 1197            \\
10 & Кобальт (Co)                                                                            & 1.02            \\
11 & Лантан (La)                                                                             & 5.99            \\
12 & Магний (Mg)                                                                             & 503             \\
13 & Марганец (Mn)                                                                           & 102             \\
14 & Мыс (Cu)                                                                                & 7.78            \\
15 & Мышьяк (As)                                                                             & 1.32            \\
16 & Натрий (Na)                                                                             & 5.51            \\
17 & Неодим (Nd)                                                                             & н/о             \\
18 & Празеодим (Pr)                                                                          & 6.59            \\
19 & Рубидий (Rb)                                                                            & 1.38            \\
20 & Самарий (Sm)                                                                            & 0.958           \\
21 & Қорғасын (Pb)                                                                           & 1.44            \\
22 & Селен (Se)                                                                              & 1.26            \\
23 & Күміс (Ag)                                                                              & 0.491           \\
24 & Скандий (Se)                                                                            & 1.14            \\
25 & Стронций (Sr)                                                                           & 59.9            \\
26 & Молбиден(Mo)                                                                            & 0.05            
\end{longtblr}

\begin{multicols}{2}
Калий гуматымен байытылған жаңа сүтқышқылды өнімге эксперименттік
зерттеу сынамалары жүргізілді. Тәжірибеде 2,5\% сүттен жасалған калий
гуматы қосылған айран, зерттеу нысаны ретінде қолданылды. Жалпы көлемі 5
литр айран дайындалды: 3 литрі тәжірибелік сынамалар үшін, 2 литрі
бақылау сынамасы. Тәжірибелік материал ретінде калий гуматының (0,25\%,
0,5\% және 1\%) ерітінділері пайдаланылды. Бақылау сынамасы 2 литр
айраннан тұрды және ешқандай қоспасыз Алматы технологиялық
университетіне (АТУ) жіберілді. Тәжірибелік айранның әр литріне калий
гуматының (0,25\%, 0,5\%, 1\%) ерітінділері қосылды.8 оқытушыдан
тұратын дегустациялық топ нәтижелерді бағалады. Нәтижелерді талдау
негізінде ең оңтайлы және қолайлы сынама ретінде 0,5\% калий гуматы
қосылған сынама таңдалды.0,1\%, 0,5\%, 1\% калий гуматының
концентраттарымен байытылған айранның органолептикалық көрсеткіштері
зерттелді. Зерттеу нәтижелері 7-ші кестеде көрсетілген.
\end{multicols}

\tcap{7-кесте. Калий гуматымен байытылған айран}
\begin{longtblr}[
  label = none,
  entry = none,
]{
  width = \linewidth,
  colspec = {Q[40]Q[196]Q[237]Q[237]Q[225]},
  cells = {c},
  cells = {font = \small},
  cell{1}{2} = {font=\bfseries\small},
  cell{1}{3} = {font=\bfseries\small},
  cell{1}{4} = {font=\bfseries\small},
  cell{1}{5} = {font=\bfseries\small},
  hlines,
  vlines,
}
№ & Көрсеткіштер                     & Қазгумат 0.1\% қосылған айран & Қазгумат 0.5\% қосылған айран & Қазгумат 1\% қосылған айран \\
1 & Органолептикалық бағалау         & I                             & II                            & III                         \\
2 & Майлылығы, \%                    & 2.5                           & 2.5                           & 2.5                         \\
3 & Қышқылдығы, \textsuperscript{0}Т & 85                            & 87                            & 90                          \\
4 & рН                               & 4,5                           & 4,43                          & 4,3                         \\
5 & Түсі                             & Ақ                            & Сүт түстес ақ                 & Ақ сұр                      
\end{longtblr}

\begin{multicols}{2}
Калий гуматымен байытылған айран зерттеулер көрсеткендей, жалпы химиялық
құрамымен ерекшеленеді, бұл оған айтарлықтай коммерциялық мүмкіндік
береді. Ашыған сүт өнімі айранды потенциалды пробиотикалық деп санауға
болады, сақтау кезінде сүтқышқылы бактериялары мен ұйытқылары болуына
байланысты оның органолептикалық көрсеткіштері жақсы деңгейде болды.

Дегустация кафедра оқытушылары үшін, сонымен қатар Көмір химиясы және
технологиясы институтының оқытушылары мен еріктілері үшін
ұйымдастырылды.

Тәжірибе. Өнімдегі ақуыздың құрамын жоғары өнімді сұйық хроматография
әдісімен анықтау.

Зерттеу нысаны дүкеннен (Адал) сатып алынған майлылығы 2,5\% айран
болды. Барлығы 1 литр айран алынды, 0,25 мл 3 түрлі ыдысқа құйылды,
бақылау үлгісіне 0,25 мл айран қалдырылды. Жұмысқа тәжірибелік материал
ретінде (0,1; 0,5 және 1)\% калий гуматы ерітінділері қолданылды.
Дегустация АТУ-нің 4 оқытушысы мен 10 студенті үшін өткізілді. Алынған
нәтижелер бойынша 0,5\% қосу арқылы ең қолайлы және ұтымды модель
таңдалды. Дайындалған сынама Алматы қаласындағы Алматы технологиялық
университетінің тамақ өнімдерінің сапасы мен қауіпсіздігін бағалау
зертханасында ақуыздың өзгеруіне зерттелді.

Зерттеу нәтижесі 0,5\% калий гуматы қосылған айранның макроэлементтік
құрамының өзгергенін, яғни бақылау айранымен салыстырғанда
макроэлементтердің мөлшерінің жоғарылағанын көрсетті. Мысалы, бақылау
айранында 0,107 мг темір және 116 мг кальций болса, 0,5\% гумат қосылған
айранда 0,122 мг темір және 157,53 мг кальций болған. Бұл нәтиже
макронутриенттердің сипатының жоғарылауын көрсетеді.

Алынған мәліметтер негізінде келесі қорытындылар жасалды: калий гуматы
негізінде айран дайындаудың рецептурасы мен технологиясы жасалды.
Әзірленген технология қарапайым, өндірісте қолдануға ыңғайлы және көп
уақытты қажет етпейді. Зерттеу барысында оңтайлы ерітінді 0,5\%
концентрациядағы калий гуматы екені анықталды. Ғылыми тәжірибелердің
нәтижелері гуминдік заттармен байытылған айран өнімдерін тәжірибе
жүзінде қолдану кезінде ешқандай жанама әсерлер немесе салдарлар
байқалмағанын көрсетеді.

Бұл бағытта әлі де клиникалық зерттеулер жалғасып жатқанын ескеріп,
денсаулыққа әсері жөнінде түбегейлі ғылыми қорытынды жасау үшін қосымша
дәлелдер қажеттігі бар екендігін атап өтеміз. Осыған байланысты
тұжырымдардың дәлелдік деңгейі нақтыланып, болашақ ғылыми зерттеу
бағыттарында анықталады.

{\bfseries Қорытынды.} Сонымен, айран өндіру технологиясында оны калий
гуматымен байыту арқылы оңтайлы нәтиже алуға болатындығы анықталды.
Осылайша гуматты қолдану арқылы айран өндіру технологиясын жетілдіру
мүмкіндігі туралы қорытынды жасалды.

Зерттеулер барысында 2,5\% майлылықтағы 1 литр айранға 0,5\% калий
гуматын қосу ұсынылды.

Калий гуматымен байытылған айранның құрамы ерекше, бұл оның коммерциялық
құндылығын арттырады. Бұл сүт өнімі сүтқышқылы бактериялары мен
ұйытқылардың болуына байланысты пробиотикалық қасиеттерге және жақсы
органолептикалық қасиеттерге ие. Айран -- өте пайдалы өнім, әсіресе
витаминдер мен микроэлементтер тапшылығы бар адамдар үшін. Сүтқышқылы
өнімдер ішек микробиомасының дұрыс тамақтануына және теңгеріміне ықпал
етеді.

Зерттеу нәтижелері бойынша негізінен гумин қышқылдарын қолдану арқылы
айран алудың рецептурасы мен технологиясы жетілдірілді. Жасалған
технологияны өндірісте пайдалану оңай: ол көп уақыт пен қаржылық
шығындарды қажет етпейді. Зерттеулер көрсеткендей, ең қолайлы ерітінді
-- 0,5\% концентрациясы бар калий гуматы. Ғылыми-экономикалық зерттеулер
және алынған нәтижелер зерттеу кезеңінде гуминді заттармен байытылған
айран өндіру кезінде теріс салдарлардың байқалмағанын көрсетеді.

Әрі қарайғы зерттеулердің болашақтағы перспективалары келесіде:

- гумин заттарының нақты фракцияларының (гумин және фульвоқышқылдары)
пробиотиктердің жекелеген штаммдарына әсерін зерттеу;

- қышқыл сүт өнімдерінің әртүрлі рецептураларына ГЗ енгізудің
стандартталған

әдістерін әзірлеу;

- ГЗ-мен байытылған өнімдердің функционалдық тиімділігін in vivo растау
үшін кең

көлемді клиникалық зерттеулер жүргізу.

Гуматты тағамдық қоспа ретінде толық пайдалану біздің ғылыми
жұмыстарымыз бойынша жалғастырылуда. Патентте (US 5626881A - «Humate
dietary supplement») адам ағзасына қоректік заттарды сіңіруді, ас
қортуды жақсарту және диетаны толықтыру мақсатында гуматтан тұратын
диеталық қоспа ұсынылған. Адам ішке тұтынуға арналған (тамақпен,
сусынмен, капсула немесе таблетка түрінде). Ғылыми-тәжірибелік модель
ретінде қолданылып жұмыстар атқарылуда.

Қорытындылай келе, калий гуматымен байытылған сүтқышқылды өнімі жаңа
қасиеттерге ие, адам денсаулығына пайдалы. Сонымен қатар экономикалық
тиімділікті төмендетпей жаңа өнімдерді дайындауға мүмкіндік бере отырып,
өзінің профилактикалық және емдік қасиеттерін сақтайды.
\end{multicols}

\begin{center}
{\bfseries Әдебиеттер}
\end{center}

\begin{refs}
1. Van Rensburg, C.E.J. The Antiinflammatory Properties of Humic
Substances: A Mini Review. //Phytotherapy Research. -2015. -Vol.29(6).
-P.791-795. DOI 101002/ptr.5319.

2. Xiong‑Xin Peng, Shuang Gai, Kui Cheng \& Fan Yang. Roles of humic
substances redox activity on environmental remediation. //\,J. Hazardous
Materials. -2022. -Vol.\,435:129070. DOI 10.1016/j.jhazmat.2022.129070.

3. Gana, W., et al. Analysis of the Impact of Selected Vitamins
Deficiencies on Health. //Nutrients. -2021, -Vol.13(9): 3163. DOI
10.3390/nu13093163.

4. Ma, Y., Cheng, X., \& Zhang, Y. The Impact of Humic Acid Fertilizers
on Crop Yield and Nitrogen Use Efficiency: A Meta-Analysis.
//Agronomy.-2024, -Vol.14(12):2763. DOI 10.3390/agronomy14122763.

5. Ampong, K., et al. Understanding the Role of Humic Acids on Crop
Production and Soil Health. //Frontiers in Agronomy. -2022. -Vol.4. -P.
1-14. DOI 10.3389/fagro.2022.848621.

6. Фасхутдинов М.\,Ф., Арынов К.\,Т., Нуркеева А.\,Б., Берикова У.,
Ошакбаев М.\,Т. Новые гуминовые удобрения из казахстанских бурых углей,
обогащённых природными биологически активными веществами: получение,
свойства и ростстимулирующая активность. //Chemical Journal of
Kazakhstan. -2020. - Т.\,4.(72)- С.113-118.

7. Anuarbekova, S., Bekshin, Z., Shaikhin, S., Alzhanova, G., Sadykov,
A., Temirkhanov, A., Sarmurzina, Z., \& Kanafin, Y.N. Exploring the
Antimicrobial and Probiotic Potential of Microorganisms Derived from
Kazakh Dairy Products. //Microbiology Research, -2024, -Vol.15(3),
-P.\,1298-1318. DOI 10.3390/microbiolres15030087.

8. Maguey Gonzalez J.\,A., Gomez‑Rosales S., Angeles M.\,L.,
Lopez‑Hernandez L.\,H., Rodriguez‑Hernandez E., Solís‑Cruz B.,
Hernandez‑Patlan D., Merino‑Guzman R., Tellez‑Isaias G. Effects of humic
acids on the recovery of different bacterial strains in an in vitro
chicken digestive model. //\,Res. Vet. Sci. -2022. -Vol.\,145.
-P.\,21-28. DOI\,10.1016/j.rvsc.2022.01.004.

9. Vaskova, J., Velika, B., Pilatova, M., Kron, I., \& Vasko, L.
Therapeutic Efficiency of Humic Acids in Intoxications. //Frontiers in
Pharmacology. -2023. -Vol.\,13(4):971 DOI 10/3390/life13040971.

10. Konysbaeva D., Yermagambet B., Kassenova Zh., Gorbulya V.,
Kazankapova M., Khaimuldinova A.Physico-Chemical Parameters of Humic
Substances of Brown Coal and Their Effect on the Growth of Phoenix
dactylifera //OnLine Journal of Biological Sciences\emph{.} -2023. -Vol.
23(4). -P.504-510. DOI 10.3844/ojbsci.2023.504.510.

11. Hudak, M., Semjon, B., Marcincakova, D., Bujnak, L., Nad, P.,
Korenekova, B., Nagy, J., Bartkovsky, M., Marcincak, S. Effect of
Broilers Chicken Diet Supplementation with Natural and Acidified Humic
Substances on Quality of Produced Breast Meat. //Animals. -2021.
-Vol.\,11(4):1087 DOI 10.3390/ani11041087.
\end{refs}

\begin{center}
{\bfseries References}
\end{center}

\begin{refs}
1. Van Rensburg, C.E.J. The Antiinflammatory Properties of Humic
Substances: A Mini Review. //Phytotherapy Research. -2015. -Vol.29(6).
-P.791-795. DOI 101002/ptr.5319.

2. Xiong‑Xin Peng, Shuang Gai, Kui Cheng \& Fan Yang. Roles of humic
substances redox activity on environmental remediation. //\,J. Hazardous
Materials. -2022. -Vol.\,435:129070. DOI 10.1016/j.jhazmat.2022.129070.

3. Gana, W., et al. Analysis of the Impact of Selected Vitamins
Deficiencies on Health. //Nutrients. -2021, -Vol.13(9): 3163. DOI
10.3390/nu13093163.

4. Ma, Y., Cheng, X., \& Zhang, Y. The Impact of Humic Acid Fertilizers
on Crop Yield and Nitrogen Use Efficiency: A Meta-Analysis.
//Agronomy.-2024, -Vol.14(12):2763. DOI 10.3390/agronomy14122763.

5. Ampong, K., et al. Understanding the Role of Humic Acids on Crop
Production and Soil Health. //Frontiers in Agronomy. -2022. -Vol.4. -P.
1-14. DOI 10.3389/fagro.2022.848621.

6. Fashutdinov M.\,F., Arynov K.\,T., Nurkeeva A.\,B., Berikova U.,
Oshakbaev M.\,T. Novye guminovye udobrenija iz kazahstanskih buryh
uglej, obogashhjonnyh prirodnymi biologicheski aktivnymi veshhestvami:
poluchenie, svojstva i roststimulirujushhaja aktivnost'.
//Chemical Journal of Kazakhstan. -2020. - T.\,4.(72)- S.113-118.{[}in
Russian{]}

7. Anuarbekova, S., Bekshin, Z., Shaikhin, S., Alzhanova, G., Sadykov,
A., Temirkhanov, A., Sarmurzina, Z., \& Kanafin, Y.N. Exploring the
Antimicrobial and Probiotic Potential of Microorganisms Derived from
Kazakh Dairy Products. //Microbiology Research, -2024, -Vol.15(3),
-P.\,1298-1318. DOI 10.3390/microbiolres15030087.

8. Maguey Gonzalez J.\,A., Gomez‑Rosales S., Angeles M.\,L.,
Lopez‑Hernandez L.\,H., Rodriguez‑Hernandez E., Solís‑Cruz B.,
Hernandez‑Patlan D., Merino‑Guzman R., Tellez‑Isaias G. Effects of humic
acids on the recovery of different bacterial strains in an in vitro
chicken digestive model. //\,Res. Vet. Sci. -2022. -Vol.\,145.
-P.\,21-28. DOI\,10.1016/j.rvsc.2022.01.004.

9. Vaskova, J., Velika, B., Pilatova, M., Kron, I., \& Vasko, L.
Therapeutic Efficiency of Humic Acids in Intoxications. //Frontiers in
Pharmacology. -2023. -Vol.\,13(4):971 DOI 10/3390/life13040971.

10. Konysbaeva D., Yermagambet B., Kassenova Zh., Gorbulya V.,
Kazankapova M., Khaimuldinova A.Physico-Chemical Parameters of Humic
Substances of Brown Coal and Their Effect on the Growth of Phoenix
dactylifera //OnLine Journal of Biological Sciences\emph{.} -2023. -Vol.
23(4). -P.504-510. DOI 10.3844/ojbsci.2023.504.510.

11. Hudak, M., Semjon, B., Marcincakova, D., Bujnak, L., Nad, P.,
Korenekova, B., Nagy, J., Bartkovsky, M., Marcincak, S. Effect of
Broilers Chicken Diet Supplementation with Natural and Acidified Humic
Substances on Quality of Produced Breast Meat. //Animals. -2021.
-Vol.\,11(4):1087 DOI 10.3390/ani11041087.
\end{refs}

\begin{info}
\hspace{1em}\emph{{\bfseries Авторлар туралы мәліметтер}}

Карденов С.А. - т.ғ.к., қауымдастырылған профессор,«С.Сейфуллин атындағы
Қазақ агротехникалық зерттеу университеті» КеАҚ, Астана,
Қазақстан,e-mail:
Askerbekovsk@mail.ru;

Игенбаев А.К. - PhD доктор, қауымдастырылған профессор, «С. Сейфуллин
атындағы Қазақ агротехникалық зерттеу университеті» КеАҚ, Астана,
Қазақстан,e-mail:aidyn\_mamyt@mail.ru;

Байтукенова Ш.Б. - т.ғ.к., қауымдастырылған профессор, «С.Сейфуллин
атындағы Қазақ агротехникалық зерттеу университеті» КеАҚ, Астана,
Қазақстан, e-mail:
baytukenova75@mail.ru;

Байтукенова С.Б. - т.ғ.к., қауымдастырылған профессор, «Қ.Құлажанов
атындағы Қазақ технология және бизнес университеті» АҚ, Астана,
Казахстан,e-mail: saule7272 @mail.ru;

Ажгереева Ж.С. - т.ғ.м., докторант, «С.Сейфуллин атындағы Қазақ
агротехникалық зерттеу университеті» КеАҚ, Казахстан, Астана,
Қазақстан,e-mail:
zhuldyz\_09.11@mail.ru.

\hspace{1em}\emph{{\bfseries Information about the authors}}

Kardenov S.A. -candidate of technical sciences, associate professor, NAO
"S.Seifullin Kazakh agrotechnical research university", Astana,
Kazakhstan,e-mail: Askerbekovsk@mail.ru;

Igenbaev A.K. -PhD doctor, associate professor, NAO"S. Seifullin Kazakh
agrotechnical research university", Astana, Kazakhstan, e-mail:
aidyn\_mamyt@mail.ru;

Baytukenova Sh.B. - candidate of technical sciences, associate
professor, NAO "S.Seifullin Kazakh Agrotechnical Research University",
Astana, Kazakhstan,e-mail: baytukenova75@mail.ru;

Baitukenova S.B. -candidate of technical sciences, associate professor,
JSC "K. Kulazhanov Kazakh University of technology and business",
Astana, Kazakhstan,e-mail: saule7272 @mail.ru;

Azhgereeva Zh.S. -candidate of technical sciences, doctoral student of
NAO "S.Seifullin Kazakh agrotechnical research university", Astana,
Kazakhstan, Kazakhstan,e-mail: zhuldyz\_09.11@mail.ru.
\end{info}
